\chapter{Методика}
\label{ch:methods}

Для того, чтобы определить вязкость воздуха в тонкой трубке, можно воспользоваться формулой Пуазейля. Для этого необходимо исследовать его ламинарное течение. В таком случае задача о течении жидкости имеет аналитическое решение.

\begin{equation}
Q  = \frac{\pi R^4 \Delta P}{8 \eta l}, 
\end{equation}
где $R$ - радиус трубки, $\Delta P$ - разность давлений на концах трубки, $\eta$ - вязкость, $l$ -длина трубки.

Это соотношение называют формулой Пуазейля. Формула Пуазейля (1) позволяет определить вязкость газа по зависимости расхода от перепада давления в трубе и используется в качестве основной расчётной формулы в данной работе. %описать что должно быть в нашей установке, чтобы измерить взякость, например можно взять трубу определённого диаметра, измерять разность давлений и расход.

Если на вход трубы поступает течение, распределение скоростей которого не является пуазейлевским, то профиль течения не может установиться сразу, а реализуется лишь на некотором расстоянии $l_\text{уст}$ от начала трубы. Оценить эту длину по порядку величины с удовлетворительной точностью можно используя выражение:

\begin{equation}
    l_\text{уст} \approx 0,2R \cdot Re, (2)
\end{equation}
где $R$ - радиус трубки, $Re$ - число Рейнольдса.

Экспериментально длину установления можно определить, измеряя распределение давления вдоль трубки $P(x)$. На неустановившемся участке будет наблюдаться отклонение от линейного закона, и при том же расходе Q градиент давления $\Delta P/l$ будет больше, чем следует из формулы Пуазейля.

Чтобы определить зависимость расхода от радиуса трубки при турбулентном течении можно принять, что флуктуации скорости в развитом турбулентном течении по порядку величины совпадают со средней скоростью потока $\Delta u \sim  \bar{u}$. При этом элементы жидкости практически равномерно перемешиваются по сечению трубы, так что в качестве «длины пробега» жидкой частицы можно взять поперечный размер системы $R$. С этими допущениями получаем, что $Q \propto R^{2.5}$ для турбулентного течения.

Для ламинарного течения это отношение можно получить из теории размерностей. Получаем, что $Q \propto R^4$.

На практике определить коэффициент $\beta$ в зависимости $Q \propto R^\beta$ можно по значениям расхода $Q$ для разных трубок при постоянном отношении разности давлений к длине трубки$\Delta P/\Delta l$ отдельно для ламинарного и турбулентного течений.